\documentclass[10pt, aspectratio=169]{beamer}
\usepackage{udec-slides}
\slide{4}{Potencial Eléctrico}

\begin{document}
\primeraDiapo

\begin{frame}{Potencial Eléctrico y Diferencia de Potencial}
    Se define como la energía potencial eléctrica por unidad de carga. La diferencia de potencial ($\Delta V$) representa el trabajo externo por unidad de carga necesario para mover una carga de prueba entre dos puntos.

    \begin{boxTeoria}{Definición Integral y Carga Puntual}
        La diferencia de potencial entre los puntos $A$ y $B$, y el potencial generado por una carga puntual $q$ a una distancia $r$ (con $V(\infty)=0$) se definen como:
        \begin{equation*}
            \Delta V = -\int_{A}^{B} \vec{E} \cdot d\vec{l} \qquad \text{y} \qquad V(r) = \frac{1}{4\pi\varepsilon_0} \frac{q}{r} \quad [\unit{\V}]
        \end{equation*}
    \end{boxTeoria}

    \begin{itemize}
        \item \textbf{Relación con el Campo:} El campo eléctrico $\vec{E}$ siempre es perpendicular a las superficies equipotenciales y apunta hacia donde el potencial disminuye.
    \end{itemize}
\end{frame}

\begin{frame}{Problema 1}

    Una carga puntual $q$ es desplazada (por un agente externo) desde el origen $(0,0)$ hasta un punto $P(L, H)$ en una región con un campo eléctrico constante $\vec{E} = E_0 \hat{x}$.

    Calcule
    \begin{itemize}
        \item La diferencia de potencial $V_{OP}$.
        \item El trabajo realizado por el agente externo.
    \end{itemize}

\end{frame}

\begin{frame}{Problema 2}

    \begin{itemize}
        \item Dos cargas idénticas $+Q$ se ubican en el eje $y$, en las posiciones $y = d$ y $y = -d$. Encuentre el potencial en un punto $P(x, 0)$ sobre el eje $x$.
        \item Un anillo de radio $R$ tiene una carga total $Q$ distribuida uniformemente ($\rho_l = \frac{Q}{2\pi R}$). Encuentre el potencial en un punto $P$ a una altura $z$ sobre su centro.
    \end{itemize}

\end{frame}

\begin{frame}{Problema 3}

    Considere dos esferas metálicas aisladas, de radios R y 3R. Si ambas esferas se encuentran al mismo potencial, ¿cuál
    es la razón entre sus cargas? Si ambas esferas poseen la misma carga, ¿cuál es la razón entre sus potenciales?

\end{frame}

\end{document}